% Root candidate, complete proof; independent written-proof audits pending.
\subsection{Actual primitive systems with permutation monodromy}

Let \(k=\overline{\mathbf F}_2\), and let \(E/k\) be the smooth
projective elliptic curve with affine equation
\[
                         y^2+y=x^3,
\]
and origin \(O\) at infinity. Put \(A=E\times E\), let
\(q:A\to E\) be the first projection, and set
\[
 V=\{O\}\times E,\qquad H=E\times\{O\},\qquad
 \Gamma=\{(a,b):a+b=O\},\qquad
 L_m=\mathcal O_A(mV+H+\Gamma)\quad(m\ge1).
                                                        \tag{EDC1}
\]
All these objects are defined over \(\mathbf F_2\). This
subsection concerns their complete linear systems, with no
identification with powers of a fixed projective embedding.

\begin{lemma}[The explicit globally generated primitive lines]
\label{lem:root-elliptic-double-cover-lines}
\status{Candidate; complete proof awaiting independent audits.}
The line \(L_m\) is ample and globally generated, of type
\((1,2m+1)\), and
\[
 L_m^2=2(2m+1),\qquad h^0(A,L_m)=2m+1,
                  \qquad L_m\cdot V=2.              \tag{EDC2}
\]
The surface \(A\) has \(p\)-rank zero.
\end{lemma}
\begin{proof}
The indicated plane cubic is smooth: on its affine chart its
partial derivative with respect to \(y\) is one, and its unique
point at infinity is nonsingular. Its invariant regular
differential is \(dx\). Cartier kills this differential because
it is exact in the function field. By Serre duality, Frobenius
on \(H^1(E,\mathcal O_E)\) is zero. Thus \(f(E)=0\), and
the product decomposition gives \(f(A)=0\).

The three elliptic divisors in (EDC1) have square zero and each
pair has intersection one. In particular
\(L_m^2=2(2m+1)\), and a fibre of \(q\) has \(L_m\)-degree
two. The line \(\mathcal O_A(mV+H)\) is a product of positive
degree elliptic lines and hence ample. The effective divisor
\(\Gamma\) is nef, as is every effective divisor on an abelian
variety; translating it away from any chosen curve proves this
intersection assertion. Therefore \(L_m\) is ample. Abelian
ample-line vanishing and Riemann--Roch give \(h^0(L_m)=2m+1\).

Relative to the product principal polarization, its polarization
homomorphism is the integral matrix
\[
                    \begin{pmatrix}m+1&1\\1&2\end{pmatrix}.
                                                        \tag{EDC3}
\]
Indeed \(V,H\) give the two diagonal entries, and \(\Gamma\)
is the pullback of \(\mathcal O_E(O)\) by addition. This
matrix has determinant \(2m+1\) and Smith invariants
\(1,2m+1\): the gcd of its entries is one. The unimodular
integer row and column operations act as automorphisms of
\(E^2\). Its kernel is therefore isomorphic to \(E[2m+1]\),
and the polarization type is \((1,2m+1)\). Since \(2m+1\)
is odd, this kernel is etale.

We give a direct proof of global generation. Set
\(L_0=\mathcal O_A(H+\Gamma)\) and
\(\mathcal V_0=q_*L_0\). Every fibre restriction of \(L_0\)
has degree two, so cohomology and base change give a vector
bundle of rank two and \(R^1q_*L_0=0\). Consequently
\[
             \deg\mathcal V_0=\chi(E,\mathcal V_0)
                 =\chi(A,L_0)=1.                    \tag{EDC4}
\]
This bundle is stable. A line subbundle \(M\subset\mathcal V_0\)
of degree \(a\ge1\) would, by adjunction, give a nonzero
section of \(L_0\otimes q^*M^{-1}\). The latter line has
self-intersection \(2-4a<0\), whereas an effective divisor
on an abelian surface is nef and has nonnegative square.
Such a subbundle is impossible. Since every proper saturated
subbundle has rank one, this proves stability of slope \(1/2\).

Projection formula gives
\(\mathcal V_m=q_*L_m=\mathcal V_0\otimes\mathcal O_E(mO)\),
which is stable of slope \(m+1/2\). For every \(b\in E(k)\),
the bundle \(\mathcal V_m^*\otimes\mathcal O_E(b)\) is stable
of negative slope \(1/2-m\), and has no nonzero section.
Indeed a nonzero section would saturate to a line subbundle of
nonnegative degree, contradicting stability. Serre duality on
\(E\) now gives
\[
                       H^1(E,\mathcal V_m(-b))=0.    \tag{EDC5}
\]
Hence \(H^0(A,L_m)\to H^0(q^{-1}(b),L_m|_{q^{-1}(b)})\)
is surjective for every \(b\). A degree-two line on an elliptic
curve is globally generated: its space of sections has dimension
two, and imposing any one point lowers this dimension to one.
Fibrewise generation and (EDC5) prove that \(L_m\) is globally
generated on \(A\).
\end{proof}

\begin{lemma}[The branch permutation module over a supersingular elliptic base]
\label{lem:root-double-cover-branch-module}
\status{Candidate; complete proof awaiting independent audits.}
Let \(\Omega\) be algebraically closed of characteristic two,
and let \(f:C\to E_\Omega\) be a finite separable morphism
of degree two from a smooth connected curve. Let \(B\) be its
reduced geometric branch set. Then \(B\ne\varnothing\), and
there is a natural isomorphism
\[
              H^1_{\mathrm{et}}(C,\mathbf F_2)
                   \simeq \mathbf F_2^B/\mathbf F_2\mathbf1.
                                                        \tag{EDC6}
\]
It is compatible with field automorphisms and with isomorphisms
of covers. In particular it is a Galois-equivariant statement
when the cover is obtained from a geometric generic fibre.
\end{lemma}
\begin{proof}
Artin--Schreier and the vanishing of Frobenius on
\(H^1(E_\Omega,\mathcal O)\) give
\[
       H^1_{\mathrm{et}}(E_\Omega,\mathbf F_2)=0,
       \qquad H^2_{\mathrm{et}}(E_\Omega,\mathbf F_2)=0.
                                                        \tag{EDC7}
\]
For the second equality, \(F-1=-1\) on this coherent
\(H^1\), and coherent \(H^2\) on the curve vanishes.
The Artin--Schreier sequence and the semilinear surjectivity
argument are recorded in
\citep[Section~59.63 and Lemma~59.63.2, complete proof]
{StacksArtinSchreierCP37}.
If \(B\) were empty, the connected degree-two etale cover
would define a nonzero class in the first group of (EDC7).
Thus the branch set is nonempty.

Write \(j:U=E_\Omega\setminus B\hookrightarrow E_\Omega\).
A separable quadratic extension in characteristic two is Galois;
let \(\iota\) be its unique nonidentity automorphism. There is
an exact sequence on the small etale site
\[
 0\longrightarrow\mathbf F_{2,E_\Omega}
  \longrightarrow f_*\mathbf F_{2,C}
  \xrightarrow{\,1+\iota\,}j_!\mathbf F_{2,U}
  \longrightarrow0.                                \tag{EDC8}
\]
Here the first map is pullback of constant functions. Off \(B\)
it is the diagonal map into \(\mathbf F_2^2\), and the last
map is the sum of the two coordinates. Over a branch point
the finite geometric fibre has one point; the first map is an
isomorphism and \(1+\iota\) is zero. More intrinsically,
\((f_*\mathbf F_2)^\iota=\mathbf F_{2,E_\Omega}\), and
the image of \(1+\iota\) inside this invariant subsheaf is
exactly \(j_!\mathbf F_{2,U}\). These stalk computations
prove (EDC8), including wild ramification.

For a finite morphism, higher direct images of any abelian
etale sheaf vanish and its direct-image stalk is the direct
sum over the finite geometric fibre. This is
\citep[Proposition~59.55.2, with Lemma~59.55.1; full proofs]
{StacksFiniteEtaleDirectCP37}. Its proof decomposes the finite
scheme over a strict henselization into strictly henselian local
schemes, on which global sections equal the closed-point stalk.
Thus it applies to \(\mathbf F_2\) in characteristic two and
requires no invertibility of the coefficient prime.

Leray, (EDC7) and the long exact sequence of (EDC8) identify
\(H^1(C,\mathbf F_2)\) with
\(H^1(E_\Omega,j_!\mathbf F_{2,U})\). The second exact
sequence
\[
 0\longrightarrow j_!\mathbf F_{2,U}
  \longrightarrow\mathbf F_{2,E_\Omega}
  \longrightarrow i_*\mathbf F_{2,B}\longrightarrow0
                                                        \tag{EDC9}
\]
for \(i:B\hookrightarrow E_\Omega\) then identifies that
group with the cokernel of the diagonal
\(\mathbf F_2\to\mathbf F_2^B\). This proves (EDC6).
Every map used is canonical for the cover; the quadratic
involution is unique. Hence these isomorphisms commute with
field automorphisms and are Galois equivariant. No finite-etale
branch scheme over an original parameter open is assumed;
the argument takes place on the geometric generic fibre and
allows inseparable residue fields of branch points.
\end{proof}

\begin{theorem}[Proper primitive monodromy in five certified prime types]
\label{thm:root-primitive-prime-permutation-counterexamples}
\status{Candidate counterexamples; complete proof awaiting independent audits.}
Let
\[
                     n\in\{5,7,11,13,17\},\qquad
                       m=(n-1)/2.
                                                        \tag{EDC10}
\]
For the explicit fixed surface and line \((A,L_m)\) in (EDC1),
the nonempty open of smooth curves with maximal relative stable
rank has genus \(n+1\), relative rank \(n-1\), and geometric
relative Lang image satisfying
\[
          G_{A,L_m}\ \subseteq\ S_n\
                     \subsetneq\operatorname{GL}_{n-1}(\mathbf F_2).
                                                        \tag{EDC11}
\]
The representation of \(S_n\) is its permutation representation
on \(\mathbf F_2^n/\mathbf F_2\mathbf1\). Consequently full
linear monodromy fails in each of these actual ample globally
generated primitive systems. In particular, the type-\((1,5)\)
example has image of order at most \(120\), whereas
\(\lvert\operatorname{GL}_4(\mathbf F_2)\rvert=20160\).
\end{theorem}
\begin{proof}
The explicit-line lemma verifies ampleness, global generation,
the precise polarization type and \(f(A)=0\). Theorem
\ref{thm:primitive-prime-generic-curve-rank} therefore applies
to each prime in (EDC10): its smooth, Euler-bijective open is
nonempty and its geometric generic curve \(C\) has
\[
                           g(C)=n+1,
                           \qquad f(C)=n-1.          \tag{EDC12}
\]
The restriction of \(q\) to \(C\) is a finite morphism of
degree \(L_m\cdot V=2\). It is separable. For example, a
purely inseparable finite morphism between these smooth curves
would be a universal homeomorphism and would preserve etale
cohomology, forcing \(f(C)=f(E)=0\), contrary to (EDC12).
Alternatively, a degree-two purely inseparable map over an
algebraically closed characteristic-two field is the relative
Frobenius up to isomorphism and preserves genus; the two genera
here differ. This uses only the function-field equality
\([k(C):k(C)^2]=2\).

Apply (EDC6) to the geometric generic double cover. Its right
side has dimension \(|B|-1\), so (EDC12) gives \(|B|=n\).
The natural Galois action on \(H^1_{\mathrm{et}}(C,\mathbf F_2)\)
therefore factors through permutations of these \(n\) branch
points. Artin--Schreier identifies this representation with
the fixed-vector space of Frobenius on coherent \(H^1(C,\mathcal O_C)\).
The ambient \(H^1(A,\mathcal O_A)\) has zero stable part.
Exactness of the stable-part functor in the actual
Frobenius-compatible hypersurface sequence identifies the latter
representation with the relative Lang representation in the
statement. Thus it is this actual relative local system, not
an unrelated cover, whose image lies in the permutation group.

For an integral normal parameter open, its fundamental group is
a quotient of the absolute Galois group of its function field;
the image of a finite local system is its generic Galois image.
The preceding geometric-generic-fibre computation consequently
proves (EDC11) on the smooth maximal-rank open. Replacing the
open by any smaller nonempty open leaves this finite image
unchanged.

The displayed permutation representation is faithful for
\(n\ge3\). If a permutation acted trivially on the quotient,
the difference of any basis vector and its image would be zero
or the all-one vector. A nonzero such difference has weight two,
and cannot be all-one when \(n\ge3\) is odd; hence every
basis vector is fixed. For \(n=5\), the group-order comparison
is \(5!=120<15\cdot14\cdot12\cdot8=20160\). The inequality
\(n!<|\operatorname{GL}_{n-1}(\mathbf F_2)|\) persists for
all larger \(n\), since increasing the dimension from \(r\)
to \(r+1\) multiplies the linear-group order by
\(2^r(2^{r+1}-1)>r+2\). This proves the strict containment
for all primes listed in (EDC10).
\end{proof}

\begin{corollary}[Disconnected nonzero Lang covers]
\label{cor:root-primitive-permutation-Lang-components}
\status{Candidate; complete proof awaiting independent audits.}
For each example in (EDC10), the nonzero relative Lang cover
has at least \((n-1)/2\) connected components. It has invariant
nonempty subcovers of total degrees
\[
                       \binom n1,\binom n2,\ldots,
                              \binom n{(n-1)/2}.     \tag{EDC13}
\]
For \(n=5\), these are subcovers of degrees five and ten.
\end{corollary}
\begin{proof}
A vector in \(\mathbf F_2^B\) is the indicator of a subset
of \(B\); passage to the quotient by constants identifies a
subset with its complement. Since \(n\) is odd, every nonzero
class has a unique representative of cardinality
\(j\in\{1,\ldots,(n-1)/2\}\). This cardinality is invariant
under every permutation of \(B\). Each of the indicated sets
is nonempty, has size \(\binom nj\), and is invariant under
the actual monodromy subgroup. The correspondence between
invariant subsets of a finite fundamental-group set and open
and closed subcovers proves (EDC13) and the component bound.
No transitivity within an individual subcover is asserted.
\end{proof}

\begin{remark}[Precise relation to the large-degree question]
\label{rem:root-primitive-permutation-scope}
\status{Candidate scope; the general A.4 question remains open.}
These are completely specified actual geometric counterexamples
to an unrestricted full-image assertion for ample globally
generated primitive systems. They show that the generic sharp
rank theorem and the birational Pfaffian theorem cannot by
themselves supply nonzero-Lang connectedness or full primitive
image. We have not identified the exact subgroup of \(S_n\)
or the theta coefficient ratio of these examples.

None of the lines \(L_m\) is very ample: its restriction to an
elliptic fibre has degree two and cannot embed that fibre.
Moreover the sequence \(L_m\) consists of different primitive
polarizations, each retaining fibre degree two. It is not the
sequence of powers of a fixed ample embedding line. Accordingly
this result does not refute Grothendieck's sufficiently-large
hypersurface-degree assertion, does not contradict any proved
fixed-principal-cube full-image case, and does not close the
general A.4 gate. Its source and elementary numerical checks are
proved in the text; no finite experiment or new numerical
certificate is used.
\end{remark}
